\documentclass[10pt,a4paper]{amsart}
\usepackage[margin=19mm]{geometry}
\usepackage{amsmath,amssymb,mathtools}
\usepackage[scr=rsfs,cal=euler]{mathalfa}
\usepackage{xfrac}
\usepackage[unicode,hidelinks,bookmarks=false]{hyperref}
\newcommand{\ie}{i.e.\ }
\newcommand{\cf}{cf.\ }
\newcommand{\resp}{respectively\ }
\newcommand{\Subsection}{\subsection}
\providecommand{\nobreakdash}{\nobreak\hbox{-}}
\setlength{\emergencystretch}{2em}
\multlinegap0pt
\usepackage{bm}
\usepackage{bbm}
\usepackage{dsfont}
\usepackage{xspace}
\usepackage{cancel}
\usepackage{mathtools}
\usepackage[shortlabels]{enumitem}
\usepackage{amsthm}
\makeatletter
\@ifundefined{theorem}{\newtheorem{theorem}{Theorem}}{}
\@ifundefined{lemma}{\newtheorem{lemma}[theorem]{Lemma}}{}
\makeatother
\newcommand{\di}{\mathrm{di}}
\newcommand{\rowsert}{\texttt{rowsert}}
\newcommand{\worsert}{\texttt{worsert}}
\title[Dyck Symmetric Functions and Applications to \(q,t\)-Catalan]{Dyck Symmetric Functions and Applications to \(q,t\)-Catalan Polynomials}
\author{Graham Hawkes}
\date{Source: arXiv:2605.13003v1 (CC BY 4.0)}
\begin{document}
\maketitle

\noindent\emph{Source and licence.} Dyck Symmetric Functions and Applications to \(q,t\)-Catalan Polynomials. Version arXiv:2605.13003v1 (CC BY 4.0), distributed under the Creative Commons Attribution 4.0 International licence, as recorded in the arXiv licence metadata for this exact version and shown on the abstract page https://arxiv.org/abs/2605.13003v1. Full rights notice: licenses/arxiv-2605.13003-CC-BY-4.0.txt.\par\smallskip

\noindent\textit{Editorial scope.} Counted unit: the Theorem whose source label is \texttt{thm:tableau-factorization-bijection} in the article (source0.tex lines 1605--1632, source label \texttt{thm:tableau-factorization-bijection}), delivered together with the article's own complete original proof of it (source0.tex lines 1633--1813, one \texttt{proof} environment of 181 lines). No mathematics is written, repaired or re-derived: statement and proof are the article's own text line for line. Required context retained uncounted: lem:rowsert-dual-dyck (lines 485--492); thm:rowsert-global-reversibility (lines 712--722); lem:rowsert-di-preservation (lines 737--747); lem:two-row-tableau-insertion (lines 1050--1071); lem:reverse-truncated-row-worsert (lines 1204--1233); lem:second-reverse-insertion-valid-two-row (lines 1418--1436). Every definition, display and label of those lines is retained unchanged. Omitted from this package and not redistributed: 1--484, 493--711, 723--736, 748--1049, 1072--1203, 1234--1417, 1437--1604, 1814--7945. Nothing inside a retained excerpt is omitted or altered: every retained body is delivered exactly as the article's lines are written, so no omission receipt of any kind accompanies this package. Citations: the retained excerpts carry no citation command at all, so no bibliography entry of the article is transcribed and no reference is redistributed. Reference adaptation: none was needed and none was made; every reference inside the delivered unit resolves inside it, so no unresolved reference or citation is produced. Printed slips kept literal: no printed word, symbol, hypothesis or number of the counted statement or of its proof was altered. Layout and class: the article's own class is \texttt{article} with packages including inputenc, fontenc, amsmath, amssymb, amsthm, amsfonts, mathtools, mathrsfs, geometry, enumitem, algorithm, algpseudocode, caption, xcolor; the wrapper is the lane's pinned \texttt{amsart} editorial wrapper at 10pt on A4, which restates the theorem-like environments the retained text needs and adds the article's own macro definitions that the retained text needs (\texttt{\textbackslash di}, \texttt{\textbackslash rowsert}, \texttt{\textbackslash worsert}), with extra wrapper packages (bm, bbm, dsfont, xspace, cancel, mathtools, enumitem). The delivered numbering is the unit's own. Assets: no retained span contains a figure environment, a graphics-inclusion command, a PDF-page inclusion, a table or a TikZ picture; no image file is copied into this package, the package declares no asset dependency and no image file is redistributed with it. Licence: version arXiv:2605.13003v1 of this article is distributed under the Creative Commons Attribution 4.0 International licence, as recorded in the arXiv licence metadata for this exact version and shown on the abstract page https://arxiv.org/abs/2605.13003v1. A full rights notice is delivered at licenses/arxiv-2605.13003-CC-BY-4.0.txt. Characters: every retained excerpt is pure ASCII, so the delivered engine, XeLaTeX with its default Latin Modern text font, reports no missing character.
\begin{lemma}[Row insertion preserves dual Dyck sequences]
\label{lem:rowsert-dual-dyck}
Let \(R_0\) and \(F_0\) be dual Dyck sequences, and let
\[
  \rowsert(R_0,F_0)=(E,R).
\]
Then both \(R\) and \(E\) are dual Dyck sequences.
\end{lemma}

\begin{theorem}[Global reversibility without Case~0]
\label{thm:rowsert-global-reversibility}
Let \(R_0\) and \(F_0\) be dual Dyck sequences.  Suppose
\[
  \rowsert(R_0,F_0)=(E,R)
\]
and that no step of this \(\rowsert\) run applies Case~0.  Then
\[
  \worsert(E,R)=(R_0,F_0).
\]
\end{theorem}

\begin{lemma}[\(\di\)-preservation for row insertion]
\label{lem:rowsert-di-preservation}
Let \(R_0\) and \(F_0\) be dual Dyck sequences, and suppose
\[
  \rowsert(R_0,F_0)=(E,R).
\]
Then
\[
  \di(R_0F_0)=\di(ER).
\]
\end{lemma}

\begin{lemma}[Two-row tableau insertion]
\label{lem:two-row-tableau-insertion}
Let \((T_0,T_1)\) be a valid two-row window in the paper top-to-bottom
convention, and let \(F\) be a dual Dyck sequence.  Apply row insertion first to
\(T_0\) and then to \(T_1\):
\[
  (G,T_0')=\rowsert(T_0,F),
  \qquad
  (E,T_1')=\rowsert(T_1,G).
\]
Assume that Case~0 never occurs in the first row insertion \(\rowsert(T_0,F)\).
Then \((T_0',T_1')\) is again a valid two-row window in the same paper
convention.  Equivalently, \(T_0'\) and \(T_1'\) are dual Dyck sequences,
\[
  |T_0'|\ge |T_1'|,
\]
and
\[
  T_0'[p]\le T_1'[p]+1
\]
for every \(0\le p<|T_1'|\).
\end{lemma}

\begin{lemma}[Reverse pass through a truncated row]
\label{lem:reverse-truncated-row-worsert}
Let \((T_0,T_1,T_2)\) be a valid three-row window in top-to-bottom row order.
Thus the rows are dual Dyck sequences,
\[
  |T_0|\ge |T_1|\ge |T_2|,
\]
and the adjacent column inequalities
\[
  T_0[p]\le T_1[p]+1,
  \qquad
  T_1[p]\le T_2[p]+1
\]
hold whenever the displayed entries exist.  Let \(F^+\) be the terminal \(k\)
entries of \(T_1\), where
\[
  0\le k\le |T_1|-|T_2|,
\]
and write
\[
  T_1=T_1^-\cdot F^+.
\]
Set
\[
  (T_1',F^-)=\worsert(T_2,T_1^-).
\]
Then Case~0 of this \(\worsert\) run is never triggered.  Moreover
\(F^-\cdot F^+\) is a dual Dyck sequence, has length at most \(|T_1|\), and
\((T_0,F^-\cdot F^+)\) is a valid two-row window in top-to-bottom row order.
\end{lemma}

\begin{lemma}[Second reverse insertion restores a two-row window]
\label{lem:second-reverse-insertion-valid-two-row}
Use the setup of Lemma~\ref{lem:reverse-truncated-row-worsert}.  Thus
\((T_0,T_1,T_2)\) is a valid three-row window in top-to-bottom row order,
\[
  T_1=T_1^-\cdot F^+,
  \qquad
  (T_1',F^-)=\worsert(T_2,T_1^-),
\]
where \(0\le k=|F^+|\le |T_1|-|T_2|\).  Put
\[
  H=F^-\cdot F^+,
\]
and apply the second reverse insertion
\[
  (T_0',F')=\worsert(H,T_0).
\]
Then \((T_0',T_1')\) is a valid two-row window in top-to-bottom row order.
\end{lemma}

\begin{theorem}[Tableau--factorization bijection]
\label{thm:tableau-factorization-bijection}
Let \(S\) be a finite multiset of integers and let \(d\ge 0\).  There is a
bijection between dual Dyck factorizations
\[
  \mathcal F=(F_0,F_1,\ldots)
\]
of the multiset \(S\) satisfying
\[
  \di(F_0F_1\cdots)=d
\]
and pairs \((P,Q)\) such that:
\begin{enumerate}[label=(\roman*)]
\item \(P\) is a Dyck tableau with entries exactly the multiset \(S\) and
      \(\di(\operatorname{RR}(P))=d\);
\item \(Q\) is a semistandard Young tableau of the same shape as \(P\);
\item the content of \(Q\) records the factorization weight, meaning
      \[
        m_i(Q)=|F_i|
      \]
      for every \(i\ge 0\).
\end{enumerate}
Equivalently, the bijection sends the factorization monomial
\[
  \prod_{i\ge 0}x_i^{|F_i|}
\]
to the semistandard-tableau weight of \(Q\).
\end{theorem}

\begin{proof}
We first prove the fixed-shape insertion step.  Fix a Young shape \(\lambda\),
an integer \(k\ge 0\), a multiset \(S\), and an integer \(d\).  We compare the
following two sets.
The first set consists of pairs \((P,F)\), where \(P\) is a Dyck tableau of
shape \(\lambda\), \(F\) is a dual Dyck sequence of length \(k\), the entries
of \(P\) together with the entries of \(F\) form the multiset \(S\), and
\[
  \di(\operatorname{RR}(P)F)=d.
\]
The second set consists of Dyck tableaux \(P'\) of shape \(\mu\), where
\(\mu\setminus\lambda\) is a horizontal strip of size \(k\), the entries of
\(P'\) form the multiset \(S\), and
\[
  \di(\operatorname{RR}(P'))=d.
\]
The local map sends \((P,F)\) to
\[
  P'=\texttt{tabsert}(P,F).
\]
Let the rows of \(P\), in the local order in which \(\texttt{tabsert}\)
processes them, be \(T_0,T_1,\ldots,T_r\), and set \(E_0=F\).  At row \(j\),
the algorithm applies
\[
  (E_{j+1},T_j')=\rowsert(T_j,E_j)
\]
and replaces \(T_j\) by \(T_j'\), continuing until the evicted sequence is empty
or until a final nonempty evicted sequence is appended as a new bottom row.  By
Lemma~\ref{lem:rowsert-dual-dyck}, every updated row and every evicted sequence
is dual Dyck.  Since each row-insertion step only moves entries among the row,
the input sequence, and the evicted sequence, the multiset of entries is
preserved.
It remains in the forward direction to check the tableau shape, the column
condition, and \(\di\).  Consider two adjacent rows during an insertion step.
If the insertion into the upper row has a terminal Case~0 phase, that phase
only appends a terminal suffix to that upper row and sends no entries farther
down.  On the preceding no-Case~0 part, Lemma~\ref{lem:two-row-tableau-insertion}
applies to the adjacent two-row window.  It shows that the updated lower row has
length at most the old upper row length and that the adjacent column condition
is preserved.  The terminal appended suffix on the upper row is to the right of
the old upper row, so it cannot create a new column violation with the lower
row.  The same argument at the bottom boundary shows that a newly appended
bottom row has length at most the old bottom row.  Therefore the final shape is
obtained from \(\lambda\) by adding a horizontal strip of size \(k\), and the
output is a Dyck tableau.
For the statistic, apply Lemma~\ref{lem:rowsert-di-preservation} at each row
step.  The invariant is most naturally stated for an augmented row-reading word.
Immediately before row \(j\) is processed, this word is
\[
  T_r\cdots T_{j+1}\,T_j\,E_j\,T_{j-1}'\cdots T_0',
\]
where \(T_r,\ldots,T_{j+1}\) are the lower rows not yet reached and
\(T_{j-1}',\ldots,T_0'\) are the already updated upper rows.  Thus the current
row step replaces the contiguous block \(T_jE_j\) by \(E_{j+1}T_j'\).  If
\(\rowsert(R,G)=(E,R')\), then
\[
  \di(RG)=\di(ER').
\]
The blocks \(RG\) and \(ER'\) have the same multiset of entries, so their
cross-\(\di\) contributions with the fixed outside context agree: such cross
terms depend only on the value counts in the moving block.  The internal
contribution is preserved by the displayed row-level equality.  Thus the
augmented word has the same \(\di\) before and after each row step.  At the end
of the process the augmented word is exactly the row-reading word of the output
tableau, including the final evicted sequence as a new bottom row if one is
created.  Hence
\[
  \di(\operatorname{RR}(P)F)=\di(\operatorname{RR}(P')).
\]
The forward local map therefore lands in the stated codomain.
We describe the inverse local map.  Let \(P'\) have shape \(\mu\) with
\(\mu\setminus\lambda\) a horizontal strip of size \(k\).  In each row, the
cells of \(\mu\setminus\lambda\) form a terminal suffix.  Write \(F_j^+\) for
that terminal suffix in row \(j\), read left-to-right, and write \(T_j^-\) for
the initial segment that remains after the suffix is removed.
Process rows from bottom to top.  Suppose the accumulated sequence passed up
from the lower part is \(E_{j+1}\).  Apply
\[
  (T_j^{\mathrm{old}},F_j^-)=\worsert(E_{j+1},T_j^-),
\]
and then pass upward the sequence
\[
  E_j=F_j^-F_j^+.
\]
At the end of the upward pass, \(E_0\) is the recovered inserted sequence
\(F\), and the recovered rows \(T_j^{\mathrm{old}}\) form the tableau \(P\).
The validity of this reverse procedure is exactly the role of
Lemmas~\ref{lem:reverse-truncated-row-worsert} and
\ref{lem:second-reverse-insertion-valid-two-row}.  The induction is on the
rows, moving upward.  The carried object \(E_{j+1}\) is not asserted to be the
row already recovered below \(j\).  Rather, it is the lower input \(T_2\) in
Lemma~\ref{lem:reverse-truncated-row-worsert} for the next three-row window.
The induction hypothesis is precisely that this carried sequence is dual Dyck
and satisfies the compatibility and length hypotheses needed to serve as that
\(T_2\).  This is true at the bottom with the empty sentinel lower row.
After the reverse pass through row \(j\), Lemma~\ref{lem:reverse-truncated-row-worsert}
gives that \(E_j=F_j^-F_j^+\) is again a dual Dyck sequence compatible with
the next row above, when such a row exists.  The adjacent recovered rows below
are validated one step later by Lemma~\ref{lem:second-reverse-insertion-valid-two-row},
when the reverse pass through the row above is performed.

At row \(j\), the number of peeled-off boxes is \(k_j=|F_j^+|\).  At the bottom
row, \(E_{j+1}\) is empty.  Otherwise, the previous lower-row pass gives
\(|E_{j+1}|\le \mu_{j+1}\).  Since \(\mu\setminus\lambda\) is a horizontal
strip, \(\mu_{j+1}\le \lambda_j=|T_j^-|\).  Hence in every row
\(|E_{j+1}|\le |T_j^-|\), which is the length hypothesis needed in
Lemma~\ref{lem:reverse-truncated-row-worsert}:
\[
  0\le k_j\le |T_1|-|T_2|.
\]
Here the lemma's \(T_1\) is the full current row \(T_j^-F_j^+\), and its
\(T_2\) is \(E_{j+1}\).  At the top row there is no further row above to check
after the pass, but the same proof as
Lemma~\ref{lem:reverse-truncated-row-worsert}, with the top row \(T_0\) and the
final column-condition conclusion omitted, still gives that Case~0 does not
occur and that \(E_0\) is dual Dyck.
Inducting upward over the rows produces a valid Dyck tableau \(P\) of shape
\(\lambda\) and a dual Dyck sequence \(F\) of length \(k\).  The entries are
preserved because each \(\worsert\) step only rearranges the current row and
accumulated sequence, and the same \(\di\)-preservation argument as above gives
\[
  \di(\operatorname{RR}(P)F)=\di(\operatorname{RR}(P'))=d.
\]
Thus the inverse local map lands in the stated domain.
The two local maps are mutual inverses.  On the no-Case~0 portions, this is the
row-level reversibility of Theorem~\ref{thm:rowsert-global-reversibility}.  The
terminal Case~0 suffixes
that are appended in the forward map are precisely the suffixes \(F_j^+\) peeled
off before applying \(\worsert\) in the inverse map.  Therefore the inverse
process reconstructs the same row states and carried sequences in reverse
order, and the fixed-shape local map is a bijection.
We iterate the local bijection.  Start with a dual Dyck factorization
\((F_0,F_1,\ldots)\) and set \(P^{(-1)}=\varnothing\).  For \(i\ge 0\), set
\[
  P^{(i)}=\texttt{tabsert}(P^{(i-1)},F_i).
\]
Only finitely many factors are nonempty.  At step \(i\), the local bijection
says that the shape grows by a horizontal strip of size \(|F_i|\).  Label the
cells added at step \(i\) by \(i\).  This produces a filling \(Q\) of the final
shape.  The filling \(Q\) is semistandard: rows are weakly increasing because
later labels are appended to row ends.  Columns are strictly increasing because
each shape difference is a horizontal strip, and the nested Young shapes force
any cell above a newly added cell in the same column to have been added at an
earlier step.  The content condition \(m_i(Q)=|F_i|\) is immediate.  Empty
factors add no cells and correspond to labels absent from \(Q\).
The local \(\di\)-preservation gives, step by step,
\[
  \di(\operatorname{RR}(P^{(i-1)})F_i)
  =
  \di(\operatorname{RR}(P^{(i)})).
\]
To pass from this local equality to the whole factorization word, include the
untouched tail in the induction.  If every factor is empty, then \(P\) is empty
and the desired \(\di\)-identity is immediate.  Otherwise, let \(N\) be the last
nonempty factor index.
We prove by induction on \(i\) that
\[
  \di(F_0F_1\cdots F_iF_{i+1}\cdots F_N)
  =
  \di(\operatorname{RR}(P^{(i)})F_{i+1}\cdots F_N),
\]
and that \(\operatorname{RR}(P^{(i)})\) has the same value multiset as
\(F_0F_1\cdots F_i\).  The multiset assertion follows at each step because
the local insertion rearranges entries.  The displayed local
\(\di\)-preservation gives equality for the processed block after adding
\(F_i\).  The remaining cross terms with the fixed tail
\(F_{i+1}\cdots F_N\) depend only on the value counts in the processed block,
so they are unchanged when \(F_0F_1\cdots F_i\) is replaced by
\(\operatorname{RR}(P^{(i)})\).  Taking \(i=N\) gives
\[
  \di(F_0F_1\cdots)=\di(\operatorname{RR}(P)),
\]
so the final tableau has parameter \(d\).
Conversely, given \((P,Q)\), let \(P^{(i)}\) be the subtableau consisting of
cells whose recording labels are at most \(i\).  Since \(Q\) is semistandard,
\(P^{(i)}\setminus P^{(i-1)}\) is a horizontal strip.  Applying the inverse
local bijection for the labels in decreasing order recovers \(P^{(i-1)}\) and
the factor \(F_i\).  Because the local maps are mutual inverses at every step,
this recovers a unique dual Dyck factorization and is inverse to the forward
construction.  The theorem follows.
\end{proof}

\end{document}
